% ===== PRÉAMBULE CANONIQUE — à coller VERBATIM en tête de CHAQUE .tex =====
\documentclass[11pt,a4paper]{article}
\usepackage[utf8]{inputenc}\usepackage[T1]{fontenc}\usepackage{lmodern}
\usepackage[french]{babel}
\usepackage{amsmath,amssymb,mathtools}\usepackage{icomma}
\usepackage[version=4]{mhchem}          % \ce{...} : équations & formules
\usepackage{siunitx}\sisetup{locale=FR} % \qty{}{}, \si{}, \ang{}
\usepackage{chemfig}                    % \chemfig{...} : structures développées
\usepackage{xcolor}\usepackage{colortbl}
\usepackage{tikz}\usepackage{pgfplots}\pgfplotsset{compat=1.18}
\usepackage[most]{tcolorbox}\usepackage{enumitem}\usepackage{array,booktabs}
\usepackage[a4paper,margin=2.2cm]{geometry}
\usetikzlibrary{arrows.meta,calc,angles,quotes,positioning,patterns,backgrounds,shapes.geometric}
\definecolor{primary}{RGB}{222,49,99}\definecolor{primarydark}{RGB}{150,22,55}
\definecolor{defcol}{RGB}{0,114,178}\definecolor{methcol}{RGB}{52,92,148}
\definecolor{excol}{RGB}{0,135,90}\definecolor{attncol}{RGB}{200,40,55}
\tcbset{boxbase/.style={enhanced,breakable,boxrule=0.9pt,arc=2.5pt,
  left=3mm,right=3mm,top=2.3mm,bottom=2.3mm,fonttitle=\bfseries,coltitle=white}}
% --- boîtes TUTORIEL (noms EXACTS, ne pas renommer) ---
\newtcolorbox{cle}[1][]{boxbase,colback=defcol!6,colframe=defcol,
  title={\textsc{À retenir}\ifx\relax#1\relax\relax\else\textmd{ : \itshape #1}\fi}}
\newtcolorbox{methode}[1][]{boxbase,colback=methcol!7,colframe=methcol,sharp corners,
  title={\textsc{Méthode}\ifx\relax#1\relax\relax\else\textmd{ --- \itshape #1}\fi}}
\newtcolorbox{exemplebox}[1][]{boxbase,colback=excol!5,colframe=excol,
  title={\textsc{Exemple}\ifx\relax#1\relax\relax\else\textmd{ : \itshape #1}\fi}}
\newtcolorbox{piege}[1][]{boxbase,colback=attncol!6,colframe=attncol,
  title={\textsc{Piège}\ifx\relax#1\relax\relax\else\textmd{ : \itshape #1}\fi}}
\newtcolorbox{remarque}[1][]{boxbase,colback=black!4,colframe=black!45,
  title={\textsc{Remarque}\ifx\relax#1\relax\relax\else\textmd{ : \itshape #1}\fi}}
% --- boîtes EXERCICE / CORRIGÉ (noms EXACTS, auto-comptées) ---
\newtcolorbox[auto counter]{exo}[1][]{boxbase,colback=primary!4,colframe=primary,
  title={\textsc{Exercice}~\thetcbcounter\ifx\relax#1\relax\relax\else\textmd{ --- \itshape #1}\fi}}
\newtcolorbox[auto counter]{corrige}[1][]{boxbase,colback=excol!4,colframe=excol,
  title={\textsc{Corrigé}~\thetcbcounter\ifx\relax#1\relax\relax\else\textmd{ --- \itshape #1}\fi}}
\newcommand{\R}{\mathbb{R}}\newcommand{\N}{\mathbb{N}}\newcommand{\Z}{\mathbb{Z}}\newcommand{\ds}{\displaystyle}
% (un \usepackage{fancyhdr}+en-tête est possible ; il est ignoré côté web)
% =========================================================================
\begin{document}
\begin{corrige}[y a-t-il résonance ?]
Il faut une double liaison conjuguée à un doublet libre (ou à une autre double).
\begin{enumerate}[label=\alph*)]
  \item \ce{CO3^2-} : \textbf{oui} (double \ce{C=O} $+$ doublets des \ce{O}).
  \item \ce{CH4} : \textbf{non} (que des liaisons simples, aucun doublet libre à délocaliser).
  \item \ce{O3} : \textbf{oui}.
  \item \ce{H2O} : \textbf{non} (aucune liaison multiple).
  \item \ce{NO3-} : \textbf{oui}.
\end{enumerate}
\end{corrige}

\begin{corrige}[combien de formes limites équivalentes ?]
Le nombre de formes $=$ nombre de positions équivalentes où placer la double liaison.
\begin{enumerate}[label=\alph*)]
  \item \ce{O3} : \textbf{2}.
  \item \ce{CO3^2-} : \textbf{3} (une par oxygène).
  \item \ce{NO3-} : \textbf{3}.
  \item \ce{NO2-} : \textbf{2}.
  \item \ce{SO2} : \textbf{2}.
\end{enumerate}
\end{corrige}

\begin{corrige}[qu'est-ce qui bouge ?]
\begin{enumerate}[label=\alph*)]
  \item \textbf{Faux} : les noyaux ne bougent jamais.
  \item \textbf{Vrai} : ce sont précisément ces électrons qui se délocalisent.
  \item \textbf{Faux} : la géométrie est identique pour toutes les formes.
  \item \textbf{Vrai} : $\leftrightarrow$ est la flèche de résonance.
\end{enumerate}
\end{corrige}

\begin{corrige}[$\leftrightarrow$ ou $\rightleftharpoons$ ?]
\begin{enumerate}[label=\alph*)]
  \item Formes limites de l'ozone : $\boldsymbol{\leftrightarrow}$ (résonance).
  \item Équilibre chimique : $\boldsymbol{\rightleftharpoons}$.
  \item Structures de résonance du carbonate : $\boldsymbol{\leftrightarrow}$.
\end{enumerate}
\end{corrige}

\begin{corrige}[conséquence sur les liaisons]
\begin{enumerate}[label=\alph*)]
  \item Elles sont \textbf{identiques} (toutes équivalentes).
  \item Leur longueur est \textbf{intermédiaire} entre celle d'une simple et d'une double liaison.
  \item La charge $-2$ est répartie \textbf{sur les trois oxygènes} ($-\tfrac{2}{3}$ chacun).
\end{enumerate}
\end{corrige}

\end{document}
